
\documentclass{article}
\usepackage{colortbl}
\usepackage{makecell}
\usepackage{multirow}
\usepackage{supertabular}

\begin{document}

\newcounter{utterance}

\centering \large Interaction Transcript for game `cladder', experiment `full\_v1.5\_default', episode 1827 with llama{-}3.1{-}8b{-}instruct.
\vspace{24pt}

{ \footnotesize  \setcounter{utterance}{1}
\setlength{\tabcolsep}{0pt}
\begin{supertabular}{c@{$\;$}|p{.15\linewidth}@{}p{.15\linewidth}p{.15\linewidth}p{.15\linewidth}p{.15\linewidth}p{.15\linewidth}}
    \# & \multicolumn{2}{c}{Player} && \multicolumn{2}{c}{Game Master} \\
    \hline

    \theutterance \stepcounter{utterance}  
    & & & \multicolumn{4}{p{0.6\linewidth}}{
        \cellcolor[rgb]{0.9,0.9,0.9}{
            \makecell[{{p{\linewidth}}}]{
                \texttt{\tiny{[P1$\langle$GM]}}
                \texttt{You are an expert in causal inference. The following question is not a typical commonsense query, but rather a meticulously designed question created by a professor specializing in causal inference, intended to assess the students' mastery of the course content.} \\
\texttt{Imagine a self{-}contained, hypothetical world with only the following conditions, and without any unmentioned factors or causal relationships: Vaccination status has a direct effect on vaccination reaction and getting smallpox. Getting smallpox has a direct effect on smallpox survival. Vaccination reaction has a direct effect on smallpox survival. Method 1: We look at how vaccination status correlates with smallpox survival case by case according to vaccination reaction. Method 2: We look directly at how vaccination status correlates with smallpox survival in general. To understand how vaccination status affects smallpox survival, is it more correct to use the Method 1 than Method 2?} \\
\texttt{Start your answer with 'Yes' or 'No', followed by additional reasoning or evidence to support your explanation.} \\
            }
        }
    }
     \\ \\

    \theutterance \stepcounter{utterance}  
    & \multicolumn{4}{p{0.6\linewidth}}{
        \cellcolor[rgb]{0.561,0.737,0.561}{
            \makecell[{{p{\linewidth}}}]{
                \texttt{\tiny{[P1$\rangle$GM]}}
                \texttt{No, it is not more correct to use Method 1 than Method 2.} \\
\\ 
\texttt{In this hypothetical world, we have a complex causal relationship between vaccination status, vaccination reaction, getting smallpox, and smallpox survival. To understand the effect of vaccination status on smallpox survival, we need to consider the potential confounding variable, which is vaccination reaction.} \\
\\ 
\texttt{Method 1, which looks at the correlation between vaccination status and smallpox survival case by case according to vaccination reaction, is actually a form of stratification or matching. By controlling for vaccination reaction, we can estimate the causal effect of vaccination status on smallpox survival within each stratum or match. This approach is often referred to as the "g{-}formula" or "causal mediation analysis."} \\
\\ 
\texttt{However, Method 2, which looks directly at the correlation between vaccination status and smallpox survival in general, is a more straightforward approach. In this case, we can use the "backdoor" or "frontdoor" adjustment to estimate the causal effect of vaccination status on smallpox survival.} \\
\\ 
\texttt{The key insight here is that the direct effect of vaccination status on smallpox survival is not the only effect we care about. We also want to understand the indirect effect of vaccination status on smallpox survival through vaccination reaction. Method 1 can help us estimate this indirect effect, but it may not provide a complete picture of the causal relationship.} \\
\\ 
\texttt{In contrast, Method 2 can provide a more comprehensive understanding of the causal relationship between vaccination status and smallpox survival, including both the direct and indirect effects. Therefore, while Method 1 can be useful for understanding the causal relationship within each stratum or match, Method 2 is a more general and comprehensive approach.} \\
\\ 
\texttt{It's worth noting that the choice between Method 1 and Method 2 ultimately depends on the research question and the goals of the analysis. If we are interested in understanding the causal relationship within each stratum or match, Method 1 may be more suitable. However, if we want to understand the overall causal relationship between vaccination status and smallpox survival, including both the direct and indirect effects, Method 2 is a more appropriate choice.} \\
            }
        }
    }
    & & \\ \\

    \theutterance \stepcounter{utterance}  
    & & & \multicolumn{2}{p{0.3\linewidth}}{
        \cellcolor[rgb]{0.9,0.9,0.9}{
            \makecell[{{p{\linewidth}}}]{
                \texttt{\tiny{[GM$|$GM]}}
                \texttt{No} \\
            }
        }
    }
    & & \\ \\

    \theutterance \stepcounter{utterance}  
    & & & \multicolumn{2}{p{0.3\linewidth}}{
        \cellcolor[rgb]{0.9,0.9,0.9}{
            \makecell[{{p{\linewidth}}}]{
                \texttt{\tiny{[GM$|$GM]}}
                \texttt{no} \\
            }
        }
    }
    & & \\ \\

    \theutterance \stepcounter{utterance}  
    & & & \multicolumn{2}{p{0.3\linewidth}}{
        \cellcolor[rgb]{0.9,0.9,0.9}{
            \makecell[{{p{\linewidth}}}]{
                \texttt{\tiny{[GM$|$GM]}}
                \texttt{game\_result = WIN} \\
            }
        }
    }
    & & \\ \\

\end{supertabular}
}

\end{document}
